\section{Qualitative Study of Solutions}
\label{sec:studio_qualitativo_edo}

It is not always possible to explicitly determine the solutions of a
differential equation. We therefore aim to find methods to determine some
salient properties of the solutions without having to determine them explicitly.
We will consider only first-order equations in normal form:
\begin{equation}\label{eqdiff}
	u'(x)= f(x,u(x))
\end{equation}
with $f\colon \Omega \subset \RR^2\to \RR$.

Recall that if $f$ satisfies the Cauchy-Lipschitz condition
(definition~\ref{def:cauchy_lipschitz}), particularly if $f$ is of class $C^1$,
then the graphs of the maximal solutions of the differential
equation~\eqref{eqdiff} \emph{fiber} the region $\Omega$ in the sense that
through every point of $\Omega$ passes a unique maximal solution
(Cauchy-Lipschitz theorem~\ref{th:cauchy_lipschitz} and
proposition~\ref{prop:separazione_soluzioni}). Furthermore,
proposition~\ref{prop:edo_massimale} tells us that the maximal solutions exit
any compact contained in $\Omega$, meaning they always reach the boundary of
$\Omega$.

\begin{example}\label{ex:edo_4630b}
Let $u$ be the maximal solution of the Cauchy problem
\[
	\begin{cases}
		u' = \cos(u^2)\cdot \sin x, \\
		u(0) = 0.
	\end{cases}
\]
Show that $u$ is defined on all of $\RR$ (thus the solution is global). Show
that $u$ is even. Show that $u$ is periodic. Show that $u(x)\ge 0$ for every
$x\in \RR$.
\end{example}
\newsavebox{\qredoquattro}\sbox{\qredoquattro}{%
\myqrshortdoclink{edo_4630b}{Qualitative study example \getrefnumber{ex:edo_4630b}}}
\begin{figure}
  \centering\includegraphics[width=\textwidth]{edo_4630b}
  \caption{Qualitative study of example~\ref{ex:edo_4630b}.
  \ifwidemargin\\\\\fi%
  \usebox{\qredoquattro}}
  \label{fig:edo_4630b}
\end{figure}
%
\begin{proof}[Solution.]
In this case, we have $f(x,y)=\cos y^2\cdot \sin x$, which is a function of
class $C^1$ on all of $\Omega=\RR^2$. Thus, the theorem of existence and
uniqueness of solutions holds. We also note (by direct verification) that
setting $H=\sqrt{\frac \pi 2}$, the constant functions $u_0(x)=-H$ and
$u_1(x)=H$ are solutions of the differential equation (\ref{eqdiff}) since
$f\enclose{x,\pm H}=0$ for every $x\in \RR$ (stationary solutions). By
proposition~\ref{prop:separazione_soluzioni}, the graph of the solution $u(x)$
cannot touch the graphs of the two functions $u_0$ and $u_1$, and since $u(0)=0$
and $u$ is a continuous function, it means that for every $x$, we have $-H <
u(x) < H$. Now consider the compact $K_M=[-M,M] \times [-H,H]$. We know that the
maximal solution $u(x)$ of the Cauchy problem under consideration must exit
every compact $K_M$ (proposition~\ref{prop:edo_massimale}). However, since
$u(x)\in (-H,H)$, it means that necessarily the graph of the maximal solution
$u(x)$ reaches the two vertical segments $x=-M$ and $x=M$, and thus the maximal
interval of existence contains the interval $[-M,M]$. Since this is true for
every $M$, the maximal interval of existence of the solution is $I=\RR$, and
thus the solution is global.

Let us prove that $u$ is even. Let $v(x)=u(-x)$. We have 
\[ 
    v'(x) = -u'(-x) 
    = -\cos(u^2(-x))\cdot \sin(-x) 
    = \cos(v^2(x))\cdot \sin(x)
\]
from which it follows that $v$ satisfies the same differential equation of which
$u$ is a solution. Moreover, $v(0)=u(0)=0$, and therefore $u$ and $v$ are
solutions of the same Cauchy problem. By the uniqueness of the solution, it must
be that $u(x)=v(x)$ for every $x$, and therefore $u(-x)=u(x)$, as we wanted to
prove.

Now let us prove that $u$ is $2\pi$-periodic. Let $w(x)=u(x+2\pi)$. Then we have
\[
  w'(x) = u'(x+2\pi) = \cos(u^2(x+2\pi))\cdot \sin(x+2\pi) 
  = \cos(w^2(x))\cdot \sin(x) 
\]
thus, again, $w$ satisfies the same differential equation of which $w$ is a
solution. Moreover, we have
\[
  u(-2\pi) - w(-2\pi) = u(2\pi) - u(0) = w(0) - u(0).
\]
This means that $u-v$ takes opposite signs at the points $a=-2\pi$ and $b=0$,
and therefore, by the intermediate value theorem~\ref{th:valori_intermedi},
there must be a point where $u-v$ vanishes. At this point, we have two solutions
of the same Cauchy problem, and thus, by the uniqueness of the solution, we must
conclude that $u(x)=v(x)$ for every $x$. This means that $u$ is $2\pi$-periodic.

Since $\abs{u(x)}<H$ with $H^2=\pi 2$, we have that $\cos(u^2(x))>0$, and
therefore $u'(x)$ has the same sign as $\sin x$: on the interval $[0,\pi]$, the
function $u(x)$ is increasing, and on $[\pi,2\pi]$, it is decreasing. Since $u$
is $2\pi$-periodic, we have that $u(2\pi)=u(0)=0$, and therefore $0$ is the
minimum value assumed by $u$.

See figure~\ref{fig:edo_4630b}.
\end{proof}

\subsection{Monotonicity and Critical Points of Solutions}

After determining the regions of $\Omega$ where the theorem of existence and
uniqueness holds, it is useful to study the sign of $f$. Indeed, the solutions
$u(x)$ of the differential equation (\ref{eqdiff}) will be strictly increasing
where $f>0$, strictly decreasing where $f<0$, and will have a critical point
where $f=0$.

A notable first case is when $f$ vanishes on a horizontal line $y=c$. In this
case, the constant function $u(x)=c$ is a solution of the differential equation
(\ref{eqdiff}). Moreover (under the hypothesis $f\in C^1$), this constant cannot
be crossed by other solutions of the differential equation.

It is important to understand whether the solutions of the differential equation
cross the curve $\{f=0\}$. Indeed, this allows us to determine the sign of the
derivative $u'(x)$ and thus obtain important information about the behavior of
the solution $u(x)$ (monotonicity, relative maxima and minima...).

More generally, it is interesting to understand when a solution of a
differential equation can cross a given curve.

\begin{theorem}\label{nonpassa}
Let $u(x)$ be a solution of the differential equation $(\ref{eqdiff})$ defined
on an interval $I$, and let $v(x)$ be any function defined on $I$, of class
$C^1$, such that its graph is entirely contained within the domain $\Omega$ of
definition of $f$. Suppose that at a fixed point $x_0$ internal to $I$, we have
$u(x_0)<v(x_0)$, and suppose further that for every $x>x_0$, we have
$v'(x)>f(x,v(x))$. Then for every $x>x_0$, we have $u(x)<v(x)$.

Similarly, if $u(x_0)>v(x_0)$ and if for every $x$, we have $v'(x)<f(x,v(x))$,
it will result that $u(x)>v(x)$ for every $x>x_0$.

Analogous results hold for $x<x_0$.
\end{theorem}

\begin{proof}
Suppose for contradiction that the set $J=\{x\in I: x\ge x_0\ \mathrm{and}\
u(x)=v(x)\}$ is not empty. This set is closed and bounded below, so if it is not
empty, it admits a minimum. Let $\bar x$ be the minimum of $J$. Certainly, $\bar
x>x_0$ since $u(\bar x)=v(\bar x)$. Moreover, for every $x\in \left[x_0,\bar
x\right[$, we have $u(x)-v(x)<0$, and therefore $u'(\bar x) \ge v'(\bar x)$
(since the left-hand incremental ratio of $u-v$ is always positive).

This is a contradiction since by hypothesis, we have $u'(\bar x) = f(\bar x,
u(\bar x)) = f(\bar x ,v(\bar x)) < v'(\bar x)$.
\end{proof}

For example, if $v(x)$ is a function whose graph is contained in $\{f=0\}$
(i.e., $f(x,v(x))=0$), then a solution $u(x)$ of the differential equation can
cross the curve $v(x)$ from above to below at points where $v'(x)>0$ and from
below to above at points where $v'(x)<0$.

\newsavebox{\qrexoqt}\sbox{\qrexoqt}{%
\myqrshortdoclink{qrexoqt}{the graph of the solution of example \getrefnumber{ex843}}}
\begin{figure}
\centering
\includegraphics[width=0.8\textwidth]{ex843.pdf}
\caption{The graph of the solution of the Cauchy problem in example~\ref{ex843}
\ifwidemargin\\\\\fi%
\usebox{\qrexoqt}}
\label{fig:ex843}
\end{figure}
%
\begin{example}\label{ex843}
Consider the following Cauchy problem
\[
\begin{cases}
  u'=1-x u^3\\
  u(0)=0
\end{cases}
\]
admits a solution $u(x)$ defined on all of $\RR$.
\end{example}
%
\begin{proof}[Solution.]
Setting $f(x,y)=1 - xy^3$, we have that $f$ vanishes on the graph of the
function $u_0(x)=1/\sqrt[3] x$. For $x>0$, the solution is increasing when it is
below the function $u_0$ and is decreasing when it is above. Moreover, since
$u_0'<0$, the solutions can cross the function $u_0(x)$ only by passing from
below to above.

Let $u(x)$ be the solution of the Cauchy problem in question, defined on a
maximal interval $I$. The initial data is $(0,0)\in \{f>0\}$. Thus, the solution
is strictly increasing in a neighborhood of $0$.

We want to show first that the solution necessarily meets both branches of the
graph of $u_0(x)$. Indeed, let $\eps>0$ be sufficiently small to belong to the
interval $I$ of existence of the solution, and let $\delta=u(\eps)$. Consider
the compact $K=\{(x,y): \eps/2 \le x \le 1/\delta^2,\ 0\le y\le
1/\sqrt[3]{x}\}$. The point $(\eps,u(\eps))$ is internal to the compact $K$, but
the solution must exit the compact for some $x>\eps$. However, the solution
cannot touch the line $y=0$ since within the compact, it remains always positive
and increasing. Nor can it touch the vertical segment $x=1/\delta^2$,
$y\in[0,\delta]$ since $u(\eps)=\delta$ and for $x>0$, the solution is strictly
increasing. Thus, necessarily the solution must meet the graph of the function
$u_0(x)=1/\sqrt[3]{x}$ at a point $\bar x>0$.

At the point $x=\bar x$, the solution $u(x)$ crosses the curve $u_0(x)$. Indeed,
at this point $u'(\bar x)=0$ while $u_0'(\bar x)<0$. Thus, in a right
neighborhood of $\bar x$, the solution is above the curve $u_0(x)$ and therefore
is decreasing (at $\bar x$, the solution presents a relative maximum).

Theorem~\ref{nonpassa} further guarantees that the solution cannot cross the
graph of the function $u_0$ again since $u'_0<0$. Thus, the solution is
decreasing and bounded. Therefore (as usual using
proposition~\ref{prop:edo_massimale}), we can deduce that the maximal solution
is defined for every $x>0$.

An analogous reasoning is done for $x<0$, obtaining the global existence of the
solution.
\end{proof}

\begin{example}
Show that the Cauchy problem
\[
\begin{cases}
	u' = -x(u^3 - \sin x) \\
	u(0) = 0
\end{cases}
\]
admits a solution $u(x)$ defined on all of $\RR$.
\end{example}
%
\begin{proof}
Let $u(x)$ be the maximal solution of the Cauchy problem defined on a maximal
interval $I$.

First, we show that $\vert u(x) \vert < 2$ for every $x\in I$. Indeed, for
$x>0$, on the curve $u_1(x)=2$, we have $f(x,2) = -x (8-\sin x) < -x < 0 =
u_1'(x)$. Thus, by Theorem~\ref{nonpassa}, the solution for $x>0$ can never
cross the line $y=2$. A similar argument is made for the line $y=-2$ and also
for $x<0$, showing that the graph of the solution can never touch the lines
$y=2$ and $y=-2$ either for $x>0$ or for $x<0$, and thus $u(x)$ is bounded (in
fact, one can be more precise and show that $\vert u(x)\vert \le 1$).

At this point, as usual, we apply Theorem~\ref{prop:edo_massimale} to the
compacts $K_M=[-M,M]\times[-2,2]$, showing that the solution must have global
existence.
\end{proof}

\subsection{Comparison Theorems}

\begin{theorem}
Let $I\subset \RR$ be an interval on which two differentiable functions $f(x)$
and $g(x)$ are defined. Let $x_1<x_2$ be two points of $I$. If $f(x_1)<g(x_1)$
and $f(x_2)>g(x_2)$, then there exists a point $\bar x \in (x_1,x_2)$ such that
$f(\bar x)=g(\bar x)$ and $f'(\bar x)\ge g'(\bar x)$.
\end{theorem}

\begin{theorem}
Consider two functions $u(x)$ and $v(x)$ solutions of the respective Cauchy
problems
\[
\begin{cases}
  u'(x)  = f(x,u(x))\\
  u(x_0) = u_0
\end{cases}
\qquad
\begin{cases}
  v'(x) = g(x,v(x))\\
  v(x_0) = v_0
\end{cases}
\]
on an interval $I$ that contains the point $x_0$. Let $\Omega\subset \RR^2$ be a
subset of the domains of definition of $f$ and $g$ such that the graphs of the
solutions $u(x)$ and $v(x)$ are contained in $\Omega$ as $x$ varies in $I$.
Suppose that for every $(x,y)\in\Omega$, we have $f(x,y)<g(x,y)$. Then for every
$x\in I$, $x> x_0$, we have $u(x)< v(x)$, while for every $x\in I$, $x<x_0$, we
have $u(x)>v(x)$.
\end{theorem}

\begin{theorem}
Let $f\colon (x_0,+\infty)\to \RR$ be a function of class $C^1$ such that the
limit
\[
  \lim_{x\to +\infty} f'(x) = \ell
\]
exists, finite or infinite. If $\ell\neq 0$, then $f$ cannot have a horizontal
asymptote as $x\to +\infty$.
\end{theorem}